% arara: clean: {extensions: [ log, pdf, toc, bbl, blg, lof, lot, toc]}
% arara: pdflatex: { shell : yes }
%
% DirectedPractice1.tex
%
% Copyright © 2019 Oromion <caznaranl@uni.pe>
%
% This program is free software: you can redistribute it and/or modify
% it under the terms of the GNU General Public License as published by
% the Free Software Foundation, either version 3 of the License, or
% (at your option) any later version.
%
% This program is distributed in the hope that it will be useful,
% but WITHOUT ANY WARRANTY; without even the implied warranty of
% MERCHANTABILITY or FITNESS FOR A PARTICULAR PURPOSE.  See the
% GNU General Public License for more details.
%
% You should have received a copy of the GNU General Public License
%s along with this program.  If not, see <http://www.gnu.org/licenses/>.

\documentclass{testfcuni} % You can find the documentation on https://github.com/carlosal1015/Notes-2019-1
%%%% Drawing's section %%%
\usepackage{tkz-fct}
\usepackage{tikz}
\usetikzlibrary{babel,calc}
\usetikzlibrary{positioning}
\usetikzlibrary{calc}

\tikzset{3D/.cd,
	x/.store in=\xx, x=0,
	y/.store in=\yy, y=0,
	z/.store in=\zz, z=0
}

\tikzdeclarecoordinatesystem{3D}{%
	\tikzset{3D/.cd,#1}%
	\pgfpoint{sin(\yy)*(\xx)}{-((\xx)/75)^2+(\zz)/100*(\xx)}%
}
\usepackage{tkz-euclide}
\usetkzobj{all}
\usetikzlibrary{decorations.markings}
\usepackage[makeroom]{cancel} % Allow cancel terms like symbolic algebra.
%%% Color's section %%%
\definecolor{bgcolor}{RGB}{247,242,222}
\definecolor{myred}{RGB}{113,22,16}
\definecolor{awesome}{rgb}{1.0, 0.13, 0.32}
\definecolor{amber}{rgb}{1.0, 0.49, 0.0}
\definecolor{azul}{RGB}{112, 44, 132}
%%% Front logo %%%
\newcommand\logo{%
	\begin{tikzpicture}
	\node[draw=lime,line width=2.5pt,fill=bgcolor,rounded corners=2pt,minimum width=0.24\textwidth,align=center,inner sep=0pt,minimum height=2.4cm] (rect)
	{ \sffamily
		\parbox{.45\linewidth}{\color{myred}\footnotesize FACULTAD\\DE CIENCIAS}%
		\parbox{.15\linewidth}{\color{myred}\Large\bfseries UNI}
	};
	\draw[violet,line width=2pt] ([xshift=18pt,yshift=-.55cm]rect.north) -- +(0,-35pt);
	\end{tikzpicture}%
}
\extrawidth{1in}

\newtcbtheorem{theorem}{Teorema}{
	fonttitle=\bfseries\upshape, fontupper=\itshape,
	colframe=green!50!black,colback=green!10!white,
	colbacktitle=green!20!white,coltitle=blue!75!black}{theo}

\begin{document}

\noindent\parbox[c]{.3\textwidth}{\logo}\hfill
\parbox[c]{1\textwidth}{\raggedright%
  \sffamily{\Large\textbf{Práctica dirigida N$\bm{^{\circ} 1}$ de Introducción al\\
      análisis convexo CM-224 Ciclo 2019-1}\footnotetext{Disponible en \url{https://github.com/carlosal1015}} \par\bigskip}
  {Undg. Carlos Alonso Aznarán Laos\\
    E-mail: \href{mailto:caznaranl@uni.pe}{caznaranl\MVAt uni.pe} 	\par\smallskip}
  {\footnotesize Science Department, National University of Engineering\par}
}

\begin{center}
  \fbox{\fbox{\parbox{6.6in}{\centering
        \textbf{Temas:} Conjuntos abiertos, cerrados.}}}
\end{center}
\hrulefill

\section*{Definiciones}

\begin{theorem}[theorem style = plain apart]{plain apart}{}%[theorem style =standard]{standard}{}
  Abcderfgh
\end{theorem}

\begin{questions}
  %\pointinleftmargin
  \pointpoints{Punto}{Puntos}
  %\bracketedpoints

  \question

  Sea ${\left(x_{k}\right)}_{k\in\mathds{N}}$ una sucesión de
  $\mathds{R}^{n}$.

  Muestre que ${\left(x_{k}\right)}_{k\in\mathds{N}}$ converge a
  $\overline{x}$ si y solo si para cualquier conjunto abierto
  $\mathcal{O}$ de $\mathds{R}^{n}$ que contiene a $\overline{x}$
  existe $\overline{k}\in\mathds{N}$ tal que para cualquier
  $k\geq\overline{k}$, los correspondientes términos $x_{k}$
  pertenece a $\mathcal{O}$.

  \begin{solution}
    \ldots
  \end{solution}

  \question

  Muestre que un conjunto $C\subset\mathds{R}^{n}$ es compacto si y
  solo si para cualquier familia de conjuntos abiertos
  ${\left\{V_{i}\right\}}_{i\in I}$ que satisface
  $C\subset\cup_{i\in I}V_{i}$ existe una subfamilia finita
  $V_{i_{1}},\ldots,V_{i_{p}}$ de $\left\{V_{i}\right\}{}_{i\in I}$
  tal que $C\subset\cup_{j=1,\ldots,p}V_{i_{j}}$.

  \begin{solution}
    \ldots
  \end{solution}

  \question

  Sea $C$ un subconjunto arbitrario de $\mathds{R}_{n}$ y
  ${\left\{V_{i}\right\}}_{i\in I}$ una familia de conjuntos abiertos
  que satisfacen $C\subset\cup_{i\in I}V_{i}$.

  Muestre que existe una subfamilia numerable
  ${\left\{V_{i_{j}}\right\}}_{j\in\mathds{N}}$ de
  ${\left\{V_{i}\right\}}_{i\in I}$ tal que
  $C\subset\cup_{j\in\mathds{N}}V_{i_{j}}$.

  \begin{solution}
    \ldots
  \end{solution}

  \question

  Dado $C\subset\mathds{R}^{n}$.

  Muestre que $\overline{x}\in\overline{C}$ si y solo si para
  cualquier $r>0$ se tiene
  $B\left(\overline{x},r\right)\cap C\neq\emptyset$.

  \begin{solution}
    \ldots
  \end{solution}

  \question

  Dado $C$ un conjunto arbitrario y $U$ un conjunto abierto, ambos de
  $\mathds{R}^{n}$.

  Asuma que $\overline{C}\cap U\neq\emptyset$.
  Muestre que $C\cap U\neq\emptyset$.

  \begin{solution}
    \ldots
  \end{solution}

  \question

  Dados $\overline{x}$ e $\overline{y}$ dos elementos diferentes de
  $\mathds{R}^{n}$.
  Tome
  \begin{math}
    \overline{z}=
    \overline{t}\overline{x}+
    \left(1-\overline{t}\right)
    \overline{y},
  \end{math}
  con $\overline{t}\in\left(0,1\right)$.

  Considere una bola abierta $B=B\left(\overline{x},r\right)$ con
  $r<\|\overline{x}-\overline{z}\|$.

  Muestre que para cualquier punto
  \begin{math}
    \overline{w}\in
    \overline{z}+
    \cup_{t>0}t
    \left(\overline{z}-B\right)
  \end{math}
  se tiene que
  \begin{math}
    z\in\overline{w}+
    \cup_{0<t\leq1}t
    \left(B-\overline{w}\right).
  \end{math}

  \begin{solution}
    \ldots
  \end{solution}

  \question

  Muestre que un subconjunto $E\subset\mathds{R}^{n}$ es cerrado si y
  solo si para cualquier sucesión
  ${\left(x_{k}\right)}_{k\in\mathds{N}}$ de $F$ que converge a algún
  punto $\overline{x}$, uno tiene $\overline{x}\in F$.

  \begin{solution}
    \ldots
  \end{solution}

  \question

  Dé un ejemplo de dos subconjuntos cerrados $C_{1}$ y $C_{2}$ de
  $\mathds{R}^{n}$ tal que la suma
  \begin{math}
    C_{1}+
    C_{2}\coloneqq
    \left\{
    x_{1}+
    x_{2}:
    x_{i}\in C_{i},
    i=1,2
    \right\}
  \end{math}
  no es cerrado.

  \begin{solution}
    \ldots
  \end{solution}

  \question
  Para una función $f\colon\mathds{R}^{m}\rightarrow\mathds{R}^{n}$,
  muestre que las dos siguientes proposiciones son equivalentes:

  \begin{parts}
    \part

    $f$ es continua,

    \part

    para cualquier conjunto abierto $V$ de $\mathds{R}^{n}$, el
    conjunto $f^{-1}\left(V\right)$ es un conjunto abierto de
    $\mathds{R}^{m}$.
  \end{parts}

  \begin{solution}
    \ldots
  \end{solution}

  \question

  Sea $f\colon\mathds{R}^{m}\rightarrow\mathds{R}^{m}$ una función
  continua e inyectiva.

  Si $U\subset\mathds{R}^{n}$ es abierto, ¿entonces $f\left(U\right)$
  es un abierto de $\mathds{R}^{n}$?

  \begin{solution}
    \ldots
  \end{solution}

  \question

  Muestre que $f\colon\mathcal{U}\rightarrow\mathds{R}^{m}$ es
  continua en $\overline{x}\in\mathcal{U}$ si y solo si para
  cualquier sucesión ${\left(x_{k}\right)}_{k\in\mathds{N}}$ de
  $\mathcal{U}$ converge a $\overline{x}$, la sucesión
  ${\left(f\left(x_{k}\right)\right)}_{k\in\mathds{N}}$ converge a
  $f\left(\overline{x}\right)$.

  \begin{solution}
    \ldots
  \end{solution}

  \question

  Asuma que $f\colon\mathds{R}^{n}\rightarrow\mathds{R}^{m}$ es una
  función continua.

  Muestre que si $W$ es un conjunto abierto de $\mathds{R}^{m}$,
  entonces su preimagen
  $f^{-1}\left(W\right)\coloneqq\left\{x\in\mathds{R}^{n}\colon f\left(x\right)\in W \right\}$
  es un conjunto abierto de $\mathds{R}^{n}$.

  ¿Qué ocurre si $W$ es un conjunto cerrado en vez de un conjunto abierto?

  \begin{solution}
    \ldots
  \end{solution}

  \question

  Sea $T\colon\mathds{R}^{p}\rightarrow\mathds{R}^{q}$ una
  transformación lineal.

  Muestre que las siguientes afirmaciones:

  \begin{parts}
    \part

    Fije una base $\mathcal{B}=\left\{v_{1},\ldots,v_{n}\right\}$ de
    $\mathds{R}^{p}$.

    Para cualquier
    $x=\left(x_{1},\ldots,x_{p}\right)\in\mathds{R}^{p}$, se
    tiene~\label{ex:13_i}
    \begin{align*}
      Tx & =\sum_{i=1}^{p}x_{i}Tv_{i} \\
         & =Ax
    \end{align*}
    donde $A$ es una matriz $q\times p$ cuyas $i$-ésima columna es
    $Tv_{i}$.

    \part

    Del item~\ref{ex:13_i} deduzca que $T$ es continua.

    \part

    Si $T$ es biyectiva, entonces $p=q$.

    Es más, $T\left(V\right)\coloneqq\left\{Tx\colon x\in V\right\}$
    es un conjunto abierto de $\mathds{R}^{q}$ si y solo si $V$ es un
    conjunto abierto de $\mathds{R}^{p}$.
  \end{parts}

  \begin{solution}
    \ldots
  \end{solution}

  \question

  Asuma que $f\colon\mathds{R}^{n}\rightarrow\mathds{R}^{m}$ es una
  función continua.

  Muestre que si $K$ es un conjunto compacto de $\mathds{R}^{m}$,
  entonces su imagen
  \begin{math}
    f\left(K\right)\coloneqq
    \left\{
    f\left(x\right)\colon
    x\in K
    \right\}
  \end{math}
  es un conjunto compacto de $\mathds{R}^{m}$.

  ¿Es la preimagen de un conjunto compacto, vía $f$, un conjunto compacto? % La función constante.

  \begin{solution}
    \ldots
  \end{solution}

  \question

  Muestre que si $S_{1}$ y $S_{2}$ son conjuntos convexos de
  $\mathds{R}^{m+n}$, entonces también lo es la suma parcial
  \[
    S=
    \left\{
    \left(x,y_{1}+y_{2}\right)\colon
    x\in\mathds{R}^{n},y_{1},y_{2}\in\mathds{R}^{m},
    \left(x,y_{1}\right)\in S_{1},
    \left(x,y_{2}\right)\in S_{2}
    \right\}.
  \]

  \begin{solution}
    \ldots
  \end{solution}

  \question

  Sea $S\subset\mathds{R}^{n}$ un conjunto.
  Asuma que $S$ es acotada.
  Muestre que
  \[
    \overline{\operatorname{co}\left(S\right)}=
    \operatorname{co}\left(\overline{S}\right).
  \]
  Mostrar mediante un ejemplo que la condición de acotamiento en
  $S$ no es superflua.

  \begin{solution}
    \ldots
  \end{solution}

  \question

  Dado $Q\in\mathds{R}^{n\times n}$.

  Muestre que si es semidefinida positiva si y solo si para cualquier
  $\lambda\in\mathds{R}$, el conjunto de subnivel
  \begin{math}
    S_{\lambda\coloneqq
    \left\{
    x\in\mathds{R}^{n}\colon
    \left\langle Qx,x\right\rangle\leq
    \lambda
    \right\}
    }
  \end{math}
  es convexa.

  \begin{solution}
    \ldots
  \end{solution}

  \question

  Dados $S$ y $C$ dos subconjuntos de $\mathds{R}^{n}$.

  Muestre que
  \[
    D\coloneqq
    \left\{
    x\colon
    x+S\subset C
    \right\}
  \]
  es un conjunto convexo si $C$ es convexo.

  \begin{solution}
    \ldots
  \end{solution}

  \question

  Sean $S_{1}$ y $S_{2}$ dos subconjuntos de $\mathds{R}^{n}$.

  Muestre que las siguientes relaciones:

  \begin{parts}
    \part

    \begin{math}
      \operatorname{co}
      \left(S_{1}\cap S_{2}\right)\subset
      \operatorname{co}\left(S_{1}\right)\cap
      \operatorname{co}\left(S_{2}\right).
    \end{math}

    Dé un ejemplo donde la inclusión estricta es cumplida.

    \part

    \begin{math}
      \operatorname{co}
      \left(S_{1}\pm S_{2}\right)=
      \operatorname{co}
      \left(S_{1}\right)\pm
      \operatorname{co}
      \left(S_{2}\right).
    \end{math}

    \part

    \begin{math}
      \operatorname{co}
      \left(S_{1}\cup S_{2}\right)=
      \operatorname{co}
      \left(
      \operatorname{co}
      \left(S_{1}\right)\cup
      \operatorname{co}
      \left(S_{2}\right)
      \right).
    \end{math}

  \end{parts}

  \begin{solution}
    \ldots
  \end{solution}

  \question

  Sean $K_{1}$ y $K_{2}$ dos conos con punta de $\mathds{R}^{n}$.

  Muestre las siguientes relaciones:

  \begin{parts}

    \part

    \begin{math}
      K_{1}+
      K_{2}=
      \operatorname{co}
      \left(
      K_{1}\cup
      K_{2}
      \right).
    \end{math}

    \part

    \begin{math}
      K_{1}\cap
      K_{2}=
      \bigcup_{\alpha\in\left[0,1\right]}
      \left(
      \alpha K_{1}\cap
      \left(1-\alpha\right)K_{2}
      \right).
    \end{math}
  \end{parts}

  \begin{solution}
    \ldots
  \end{solution}

  \question
  Sea $X\subset\mathds{R}^{m}$.

  Muestre que las siguientes relaciones:

  \begin{parts}

    \part

    $X$ y el $\operatorname{co}$ tienen la misma cápsula afín.

    \part

    \begin{math}
      \operatorname{cono}
      \left(X\right)=
      \operatorname{cono}
      \left(
      \operatorname{co}
      \left(X\right)
      \right).
    \end{math}

    \part

    \begin{math}
      \operatorname{aff}
      \left(
      \operatorname{co}
      \left(X\right)
      \right)\subset
      \operatorname{aff}
      \left(
      \operatorname{cono}\left(X\right)
      \right).
    \end{math}

    Dé un ejemplo con la inclusión estricta.

    \part

    Si $0\in\operatorname{co}\left(X\right)$, entonces
    \begin{math}
      \operatorname{aff}
      \left(
      \operatorname{co}
      \left(X\right)
      \right)=
      \operatorname{aff}
      \left(
      \operatorname{cono}
      \left(X\right)
      \right).
    \end{math}
  \end{parts}

  \begin{solution}
    \ldots
  \end{solution}

  \question

  Sean $d_{1},d_{2},\ldots,d_{p}$ una cantidad finita de elementos
  de $\mathds{R}^{n}$.

  Muestre que
  \[
    C=
    \left\{
    x\in\mathds{R}^{n}\colon
    x=
    \sum_{i=1}^{p}
    \lambda_{i}d_{i},\quad
    \lambda_{1},\ldots,\lambda_{p}\in
    \left[0,\infty\right)
    \right\}
  \]
  es un cono convexo cerrado.

  \begin{solution}
    \ldots
  \end{solution}

  \question
  Sea $C\subset\mathds{R}^{n}$ y $\lambda_{1},\lambda_{2}\in\left[0,\infty\right)$.

  Muestre que si $C$ es un conjunto convexo, entonces
  \[
    \left(
    \lambda_{1}+
    \lambda_{2}
    \right)
    C=
    \lambda_{1}C+
    \lambda_{2}C.
  \]
  Muestre mediante un ejemplo que la condición de convexidad en $C$ no es superflua.

  \begin{solution}
    \ldots
  \end{solution}

  \question

  Sea $T\colon\mathds{R}^{p}\to\mathds{R}^{q}$ una transformación lineal.

  Si $C\subset\mathds{R}^{p}$ es un conjunto convexo, entonces también lo es la imagen $T\left(C\right)$.

  ¿Se cumple la implicación recíproca o conversa?

  \begin{solution}
    \ldots
  \end{solution}

  \question

  Sea $T\colon\mathds{R}^{p}\rightarrow\mathds{R}^{q}$ una transformación lineal y $C\subset\mathds{R}^{p}$.

  Entonces
  \[
    T\left(\operatorname{co}\left(C\right)\right)=
    \operatorname{co}\left(T\left(C\right)\right)\quad
    \text{and}\quad
    T\left(\operatorname{aff}\left(C\right)\right)=
    \operatorname{aff}\left(T\left(C\right)\right).
  \]

  \begin{solution}
    \ldots
  \end{solution}

  \question

  Sea $T\colon\mathds{R}^{p}\rightarrow\mathds{R}^{q}$ una transformación lineal.
  Si $C\subset\mathds{R}^{p}$ es un conjunto convexo, entonces
  \[
    T\left(\operatorname{ri}\left(C\right)\right)=
    \operatorname{ri}\left(T\left(C\right)\right).
  \]
  ¿También se mantiene con la clausura en lugar del interior
  relativo?

  \begin{solution}
    \ldots
  \end{solution}

  \question

  Sea $C_{1}$ y $C_{2}$ dos conjuntos convexos de $\mathds{R}^{n}$ que satisfacen $C_{1}\subset C_{2}$.

  Muestre que si
  \begin{math}
    \operatorname{aff}
    \left(C_{1}\right)=
    \operatorname{aff}
    \left(C_{2}\right),
  \end{math}
  entonces
  \begin{math}
    \operatorname{ri}
    \left(C_{1}\right)\subset
    \operatorname{ri}
    \left(C_{2}\right).
  \end{math}

  \begin{solution}
    \ldots
  \end{solution}

  \question

  Para una cantidad de $n\times n$ matrices simétricas
  $A_{0},A_{1},\ldots,A_{m}$, considere para cada
  $x=\left(x_{1},\ldots,x_{m}\right)\in\mathds{R}^{m}$,
  \[
    A\left(x\right)=
    A_{0}+
    \sum_{i=1}^{m}
    x_{i}A_{i}.
  \]

  Muestre que el conjunto
  \begin{math}
    C=
    \left\{
    x\in\mathds{R}^{m}\colon
    A\left(x\right)
    \text{ es semidefinida positiva}.
    \right\}
  \end{math}
  es un conjunto cerrado y convexo.

  ¿Qué es la cápsula afín de esta cápsula cónica?

  \begin{solution}
    \ldots
  \end{solution}

  \question Sea

  \begin{solution}
    \ldots
  \end{solution}

\end{questions}

\begin{flushright}
  Facultad de Ciencias, \today.
\end{flushright}
\end{document}
